\section{A signed kernel for actual square differences}
\label{sec:square-kernel}

The preceding tuple measures impose square-difference constraints in residue fields.
This section supplies the passage to actual positive integer squares.
We construct a signed measure supported on such squares and approximate
its Fourier transform uniformly by a sum over small rational frequencies.
Because every actual square shift vanishes between the supports under
consideration, the signs of the measure cause no difficulty. The weighted
square sums and major/minor-arc analysis have a direct antecedent
in~\cite[Section~5]{GreenSawhney2025}. Here the signed divisor weight
is chosen to match the strict-square pair law of our finite-field tuples.

For the remainder of the paper, fix
\begin{equation}
 \beta=10^{-5}.
 \label{eq:kernel-beta}
\end{equation}
Write $e(t)=\exp(2\pi i t)$, and let $q(\xi)$ be the reduced denominator
of $\xi\in\mathbb Q/\mathbb Z$, with $q(0)=1$.
For every rational frequency define the normalized unit quadratic sum
\begin{equation}
 \mathfrak g(\xi)
 =\frac{1}{\varphi(q(\xi))}
   \sum_{m\in(\mathbb Z/q(\xi)\mathbb Z)^\times}e(\xi m^2),
 \qquad \mathfrak g(0)=1.
 \label{eq:kernel-unit-gauss}
\end{equation}
Here $\varphi$ is Euler's totient function.
If $D=8d_0$ with $d_0$ odd and squarefree, call a residue modulo $D$
a \emph{strict square class} if it is $1$ modulo $8$ and a nonzero
square modulo every prime dividing $d_0$.

\begin{proposition}[Ordered square-kernel estimate]
\label{prop:square-kernel}
Let $N\ge1$ be an integer, put $H=N^{\beta/2}$, and let
$D=8d_0\le N^\beta$, where $d_0$ is odd and squarefree.
Suppose $F,J:[N]\to\mathbb C$ satisfy $|F|,|J|\le1$ and are supported
in residue classes $a,b$ modulo $D$, respectively. Assume that $b-a$
is a strict square class, that $x<y$ whenever
$x\in\operatorname{supp}F$ and $y\in\operatorname{supp}J$, and that
no such difference $y-x$ is the square of a positive integer.
With
\[
 \widehat F(\xi)=\frac1N\sum_{n=1}^N F(n)e(-n\xi),
 \qquad
 \widehat J(\xi)=\frac1N\sum_{n=1}^N J(n)e(-n\xi),
\]
one has
\begin{equation}
 D\left|
 \sum_{\substack{\xi\in\mathbb Q/\mathbb Z:\ q(\xi)\le H\\
                  q(\xi)\ \mathrm{squarefree},\ (q(\xi),D)=1}}
      \mathfrak g(\xi)\widehat F(-\xi)\widehat J(\xi)
 \right|
 \ll H^{-1/3}.
 \label{eq:kernel-conclusion}
\end{equation}
The implied constant is absolute and is independent of all the data
in the statement.
\end{proposition}

\begin{proof}
Put $r=b-a$ modulo $D$, and denote the frequency set in
\eqref{eq:kernel-conclusion} by $\Xi$.
We first define the model and actual-square kernels. We then compare
their Fourier transforms on minor and major arcs, and finish with the
exact bilinear Fourier identity.

\paragraph{The two kernels.}
For every odd prime $p\mid d_0$, choose one root $c_p$ of $r$ modulo $p$.
The conditions $m\equiv c_p\pmod p$ for all such primes, together
with $m$ odd, specify one residue class modulo
$D'=2d_0=D/4$. Let $\rho_D$ be $D'$ times the indicator of this class.
Thus $\rho_D$ has periodic mean one and supremum $D'$.
For every odd prime $p\nmid D$, set
\begin{equation}
 w_p(m)=\frac{p\mathbf1_{p\mid m}-1}{p-1}
       =\frac1{p-1}\sum_{j=1}^{p-1}e(jm/p),
 \qquad
 w_u(m)=\prod_{p\mid u}w_p(m),
 \label{eq:kernel-sieve-weight}
\end{equation}
where $u$ is squarefree and coprime to $D$, and $w_1=1$.
The second equality follows by summing all characters modulo $p$.
In particular $|w_u|\le1$.

The factor $1-w_p(m)=p(p-1)^{-1}1_{p\nmid m}$ is the density of
the uniform unit law modulo $p$ relative to the uniform law. Hence for
any positive integer $q$ coprime to $D$,
\begin{equation}
 \sum_{u\mid\operatorname{rad}(q)}(-1)^{\omega(u)}w_u(m)
 =\prod_{p\mid q}(1-w_p(m))
 =\frac q{\varphi(q)}1_{(m,q)=1}.
 \label{eq:kernel-complete-sieve}
\end{equation}
Here $\omega(u)$ counts the distinct prime divisors of $u$, and
$\operatorname{rad}(q)$ is the product of the distinct primes dividing $q$. Thus the complete divisor sum
converts a uniform residue into a uniform unit, whose square has Fourier
multiplier $\mathfrak g$. We truncate this sum at $u\le H$ in the
actual-square kernel. At a rational frequency with denominator $q\le H$
coprime to $D$, all divisors of $\operatorname{rad}(q)$ remain present;
the periodic mean of every term containing a prime outside $q$ vanishes. This will give the
exact low-denominator comparison. The other rational coefficients will
be bounded by quadratic Gauss sums.

Define two finitely supported functions on $\mathbb Z$ by
\begin{align}
 k_0(n)
 &=\frac DN\mathbf1_{1\le n\le N}\mathbf1_{n\equiv r\ (D)}
       \sum_{\xi\in\Xi}\mathfrak g(\xi)e(-\xi n),
       \label{eq:kernel-model}\\
 k(n)
 &=\sum_{1\le m\le\sqrt N}\mathbf1_{n=m^2}\frac{2m}{N}\rho_D(m)
       \sum_{\substack{u\le H\ \mathrm{squarefree}\\(u,D)=1}}
                   (-1)^{\omega(u)}w_u(m).
       \label{eq:kernel-actual}
\end{align}
The first kernel represents the rational expression in the proposition;
the second is supported on actual positive integer squares. The factor
$2m/N$ is the discrete counterpart of the substitution
$dn=2m\,dm$, matching the uniform weight $1/N$ in the model kernel.
For these kernels use the positive-sign Fourier transforms
\[
 S_0(\alpha)=\sum_{n\in\mathbb Z}k_0(n)e(\alpha n),
 \qquad
 S(\alpha)=\sum_{n\in\mathbb Z}k(n)e(\alpha n),
 \qquad \alpha\in\mathbb R/\mathbb Z.
\]
Our immediate goal is
\begin{equation}
 \sup_{\alpha\in\mathbb R/\mathbb Z}|S(\alpha)-S_0(\alpha)|
 \ll H^{-1/3}.
 \label{eq:kernel-uniform-approximation}
\end{equation}

Put $V_N(t)=N^{-1}\sum_{n=1}^N e(tn)$.
Expanding the congruence indicator in \eqref{eq:kernel-model} gives
\begin{equation}
 S_0(\alpha)=\sum_{\xi\in\Xi}\sum_{j\bmod D}
    \mathfrak g(\xi)e(-jr/D)V_N(\alpha-\xi+j/D).
 \label{eq:kernel-model-centers}
\end{equation}
There are at most $DH^2$ terms, each of coefficient modulus at most
one, and their centers have denominators at most $DH$.
The centers are distinct. Indeed, an equality
$\xi-j/D=\xi'-j'/D$ would make $\xi-\xi'$ have denominator both
coprime to $D$ and dividing $D$, so that $\xi=\xi'$ and $j=j'$ modulo $D$.

Every $\alpha_0\in\mathbb Q/\mathbb Z$ has a unique decomposition
$\alpha_0=\alpha_D+\alpha_1$ in which the denominator of $\alpha_D$
has only prime factors dividing $D$ and $q_1=q(\alpha_1)$ is coprime
to $D$. Thus the coefficient of the center $\alpha_0$ in
\eqref{eq:kernel-model-centers}, taken to be zero when it is absent, is
\begin{equation}
 c_0(\alpha_0)=
 \begin{cases}
 e(\alpha_Dr)\mathfrak g(\alpha_1),&
 q(\alpha_D)\mid D,\quad q_1\le H,\quad q_1\text{ squarefree},\\
 0,&\text{otherwise}.
 \end{cases}
 \label{eq:kernel-model-coefficient}
\end{equation}
The sign is worth recording: in \eqref{eq:kernel-model-centers},
$\alpha_D=-j/D$, and its coefficient contains $e(-jr/D)$.

\paragraph{Minor arcs.}
We initially assume $N$ is sufficiently large by an absolute amount.
Take the major arcs to be the closed balls of radius $2N^{-9/10}$
about all reduced rationals of denominator at most $N^{1/10}$.
They are disjoint for sufficiently large $N$, since two distinct
centers are at least $N^{-1/5}$ apart. The minor arcs are their complement.
We write $\|t\|$ for the distance of a real number $t$ to the nearest integer.

For $\alpha$ on the minor arcs, Dirichlet approximation with
$Q_0=\lfloor N^{9/10}\rfloor$ gives coprime integers $a,q$ satisfying
\[
 1\le q\le Q_0,
 \qquad |\alpha-a/q|\le\frac1{qQ_0}\le\frac1{q^2},
\]
where the difference is represented by a real number of least absolute value.
For completeness, partitioning the circle into $Q_0$ intervals and
placing the $Q_0+1$ points $0,\alpha,\ldots,Q_0\alpha$ in them gives
integers $1\le h\le Q_0$ and $b$ with $|h\alpha-b|\le Q_0^{-1}$.
Reducing $b/h$ to $a/q$ gives the displayed bounds.
If $q\le N^{1/10}$, the first approximation bound would put $\alpha$
within $2N^{-9/10}$ of a major-arc center. Consequently
\begin{equation}
 N^{1/10}<q\le N^{9/10}.
 \label{eq:kernel-minor-denominator}
\end{equation}

Let $M=\lfloor\sqrt N\rfloor$. Uniformly in real $\theta$ and
integers $1\le L\le M$, differencing and a geometric sum give
\begin{equation}
 \left|\sum_{m=1}^L e(\alpha m^2+\theta m)\right|^2
 \ll M+\sum_{j=1}^M\min\bigl(M,\|2\alpha j\|^{-1}\bigr).
 \label{eq:kernel-differencing}
\end{equation}
Indeed, pairs at distance $j$ leave a linear phase with coefficient
$2\alpha j$ in the inner sum. Its remaining factors, including
$e(\theta j)$, have modulus one. This also proves the asserted
uniformity in the linear twist.

Partition the $j$ indices into consecutive blocks of length
$\lfloor q/4\rfloor$. Two distinct indices in a block differ by
$0<h<q/4$. The residue $2ah$ is nonzero modulo $q$: since $(a,q)=1$,
its vanishing would force $q\mid2h$, whereas $0<2h<q/2$.
This argument includes even $q$. Hence
\[
 \|2\alpha h\|\ge\|2ah/q\|-2h/q^2\ge\frac1{2q}.
\]
In each block the points $2\alpha j$ are therefore separated on the
circle by at least $1/(2q)$. Ordering their distances to zero bounds
the contribution of one block by
$O(M+q\log(2q))$: only a bounded number can lie within distance
$1/q$ of zero, and the others contribute a harmonic sum.
There are $O(M/q+1)$ blocks, so
\begin{equation}
 \sum_{j=1}^M\min\bigl(M,\|2\alpha j\|^{-1}\bigr)
 \ll (M/q+1)\bigl(M+q\log(2q)\bigr).
 \label{eq:kernel-spacing-sum}
\end{equation}
By \eqref{eq:kernel-minor-denominator},
\eqref{eq:kernel-differencing}--\eqref{eq:kernel-spacing-sum} imply
\[
 \sup_{\theta\in\mathbb R}\ \max_{L\le M}
 \left|\sum_{m=1}^L e(\alpha m^2+\theta m)\right|
 \ll N^{9/20}\sqrt{\log N}.
\]
Partial summation with the weight $2m/N$ gives
\begin{equation}
 \sup_{\theta\in\mathbb R}
 \left|\sum_{m\le\sqrt N}\frac{2m}{N}
                    e(\alpha m^2+\theta m)\right|
 \ll N^{-1/20}\sqrt{\log N}.
 \label{eq:kernel-minor-quadratic}
\end{equation}

For a periodic function, use Fourier coefficients normalized by its
period. The Fourier coefficient $\ell^1$ norm of $\rho_D$ is $D'$:
its character expansion has $D'$ coefficients of modulus one.
By \eqref{eq:kernel-sieve-weight}, the corresponding norm of each
$w_p$ is one, and that of $w_u$ is at most one by the convolution
inequality. The entire periodic weight multiplying the quadratic
phase in $S$ thus has Fourier coefficient $\ell^1$ norm at most
$D'\lfloor H\rfloor\le DH$. Expansion into these characters and
\eqref{eq:kernel-minor-quadratic} yield
\begin{equation}
 |S(\alpha)|\ll DH N^{-1/20}\sqrt{\log N}
 \qquad(\alpha\text{ on the minor arcs}).
 \label{eq:kernel-minor-actual}
\end{equation}

Every center in \eqref{eq:kernel-model-centers} has denominator at
most $DH\le N^{3\beta/2}<N^{1/10}$, so it is a major-arc center.
On the minor arcs its distance from $\alpha$ exceeds $2N^{-9/10}$.
The elementary estimate
\[
 |V_N(t)|\le\min\bigl(1,(2N\|t\|)^{-1}\bigr)
\]
therefore gives
\begin{equation}
 |S_0(\alpha)|\ll DH^2N^{-1/10}
 \qquad(\alpha\text{ on the minor arcs}).
 \label{eq:kernel-minor-model}
\end{equation}

\paragraph{Reduction on a major arc to periodic means.}
Write $\alpha=\alpha_0+y$, where
$q=q(\alpha_0)\le N^{1/10}$ and $|y|\le2N^{-9/10}$.
If $\alpha_c\ne\alpha_0$ is a center in
\eqref{eq:kernel-model-centers}, rational separation gives
$\|\alpha_c-\alpha_0\|\ge1/(qDH)$.
Because $qDH\le N^{1/10+3\beta/2}$, for all sufficiently large $N$
we have $\|\alpha-\alpha_c\|\ge1/(2qDH)$.
The contributions of these other centers sum to
\begin{equation}
 S_0(\alpha)=c_0(\alpha_0)V_N(y)+O(qD^2H^3/N).
 \label{eq:kernel-major-model}
\end{equation}

For an admissible $u$, put
\[
 a_u(m)=\rho_D(m)w_u(m)e(\alpha_0m^2),
 \qquad
 A_u=\frac1{T_u}\sum_{m\bmod T_u}a_u(m),
 \qquad T_u=\operatorname{lcm}(q,D',u).
\]
The sequence $a_u$ has period $T_u\le qDu$ and supremum at most $D$.
Breaking a partial sum into complete periods and one remaining interval
shows, for every real $t\ge0$, that
\[
 \left|\sum_{1\le m\le t}a_u(m)-A_ut\right|\ll DT_u.
\]
This includes the fractional part of $t$.
For $f_y(x)=(2x/N)e(yx^2)$ on $[0,\sqrt N]$, direct differentiation gives
\[
 |f_y(\sqrt N)|+\int_0^{\sqrt N}|f_y'(x)|\,dx
 \ll\frac{1+|y|N}{\sqrt N}.
\]
Stieltjes partial summation consequently gives
\begin{align*}
 \sum_{m\le\sqrt N}\frac{2m}{N}e(ym^2)a_u(m)
 &=A_u\int_0^{\sqrt N}\frac{2x}{N}e(yx^2)\,dx\\
 &\hspace{5mm}+O\!\left(\frac{qD^2u(1+|y|N)}{\sqrt N}\right).
\end{align*}
The integral equals $N^{-1}\int_0^N e(yt)\,dt$.
Comparing the integral to its integer Riemann sum bounds its difference
from $V_N(y)$ by $O((1+|y|N)/N)$. Since $|A_u|\le D$, that error
is dominated by the preceding displayed error.
Summing over $u\le H$ proves
\begin{equation}
 S(\alpha)=c(\alpha_0)V_N(y)
       +O\!\left(\frac{qD^2H^2(1+|y|N)}{\sqrt N}\right),
 \qquad
 c(\alpha_0)=\sum_{\substack{u\le H\ \mathrm{squarefree}\\(u,D)=1}}
                  (-1)^{\omega(u)}A_u.
 \label{eq:kernel-major-actual}
\end{equation}
The use of $\sqrt N$ as a real endpoint in this calculation accounts
for rounding, while the sum itself contains only integers $m$ with
$m^2\le N$.

We have reduced the major-arc comparison to $c(\alpha_0)$ and
$c_0(\alpha_0)$. We next compute the contribution at primes dividing
$D$, and then the divisor sieve on the coprime denominator $q_1$.

\paragraph{The prescribed roots, including the prime two.}
At an odd prime $p\mid D$, the local factor of $\rho_D$ is
$p\mathbf1_{m\equiv c_p\ (p)}$.
For every $j\ge1$, squaring maps the lifts of $c_p$ modulo $p^j$
bijectively to the residues congruent to $r$ modulo $p$.
Indeed, if $x$ is such a lift, then
\[
 (x+tp^k)^2\equiv x^2+2xtp^k\pmod {p^{k+1}},
\]
and $2x$ is invertible modulo $p$. This proves the bijection
inductively over $k$.
For a primitive additive character $e(am^2/p^j)$, its expectation
under this normalized root restriction is therefore $e(ar/p)$ if
$j=1$, and zero if $j\ge2$: in the latter case one averages a
nontrivial character over all the lifts of a single residue modulo $p$.
If no $p$-part is present in the character, the expectation is one.

At two the local factor of $\rho_D$ is $2\mathbf1_{m\text{ odd}}$.
For $j\le3$, every odd square is $r$ modulo $2^j$, so the expectation
of $e(am^2/2^j)$ is $e(ar/2^j)$.
For $j\ge3$, the kernel of the squaring homomorphism on
$(\mathbb Z/2^j\mathbb Z)^\times$ has exactly four elements.
To see this, if $(x-1)(x+1)$ is divisible by $2^j$, exactly one of
these two consecutive even factors has 2-adic valuation one, and
the other is divisible by $2^{j-1}$.
The four solutions modulo $2^j$ are
$1,-1,1+2^{j-1},-1+2^{j-1}$.
The image of squaring thus has $2^{j-3}$ elements.
Every element of this image is $1$ modulo $8$, and this residue class
also has $2^{j-3}$ elements. The square of a uniform odd residue
modulo $2^j$ is consequently uniform among all residues equal to
$1$ modulo $8$. For $j>3$, a primitive character averages to zero
on that class.

Apply these local facts to the decomposition
$\alpha_0=\alpha_D+\alpha_1$ already used in
\eqref{eq:kernel-model-coefficient}. Chinese remaindering shows that
the factor at primes dividing $D$ is
\begin{equation}
 \mathbb E_m\rho_D(m)e(\alpha_Dm^2)
 =\begin{cases}
   e(\alpha_Dr),&q(\alpha_D)\mid D,\\
   0,&q(\alpha_D)\nmid D.
  \end{cases}
 \label{eq:kernel-root-law}
\end{equation}
The expectation here is the uniform average over any common period.
The normalized odd-root restrictions, and the normalized oddness
restriction at two, are precisely responsible for the factor $D$
in the model kernel. In particular, choosing only one odd root
introduces no additional factor in \eqref{eq:kernel-root-law}.

\paragraph{The coprime divisor sieve.}
Every prime of an admissible $u$ is coprime to $D$.
If such a prime does not divide $q_1$, its mean-zero factor $w_p$
is independent of the exponential and all other factors by the
Chinese remainder theorem. Thus $A_u=0$ unless
$u\mid\operatorname{rad}(q_1)$.

Suppose first that $q_1\le H$. All divisors of
$\operatorname{rad}(q_1)$ then occur in the sum. By
\eqref{eq:kernel-complete-sieve}, the divisor weight is the density of
the uniform unit law modulo $q_1$ relative to the uniform law.
If $p^j$ exactly divides $q_1$ with
$j\ge2$, its local primitive unit quadratic average vanishes.
Indeed, conditional on a unit $x$ modulo $p^{j-1}$, the $p$ lifts
$x+tp^{j-1}$ have phase factors differing by $e(2atx/p)$,
where $p\nmid 2ax$. Their sum is zero.
If $q_1$ is squarefree, the unit average is exactly
$\mathfrak g(\alpha_1)$.
Together with \eqref{eq:kernel-root-law}, this proves
\begin{equation}
 c(\alpha_0)=c_0(\alpha_0)\qquad(q_1\le H).
 \label{eq:kernel-small-coefficient}
\end{equation}

For $q_1>H$ the model coefficient is zero, and we need a bound for
the truncated divisor sum. We give the prime-power calculation
explicitly. For an odd prime $p$, an integer $j\ge1$, and $p\nmid a$,
write
\[
 G_j(a)=\frac1{p^j}\sum_{m\bmod p^j}e(am^2/p^j),
 \qquad G_0(a)=1.
\]
Because two is invertible modulo $p^j$, the substitution
$(m,n)\mapsto(m-n,m+n)$ is a bijection. Consequently
\begin{equation}
 |G_j(a)|^2
 =\frac1{p^{2j}}\sum_{u,v\bmod p^j}e(auv/p^j)
 =p^{-j}.
 \label{eq:kernel-prime-power-gauss}
\end{equation}
Let
$T_j=\mathbb E_{m\bmod p^j}\mathbf1_{p\mid m}e(am^2/p^j)$.
For $j=1$ one has $T_1=1/p$; for $j\ge2$, substituting $m=pt$
gives $T_j=p^{-1}G_{j-2}(a)$.
It follows that
\begin{equation}
 B_j:=\mathbb E_{m\bmod p^j}w_p(m)e(am^2/p^j)
      =\frac{pT_j-G_j(a)}{p-1},
 \qquad |B_j|\le2p^{-j/2}.
 \label{eq:kernel-weighted-prime-power-gauss}
\end{equation}
For $j\ge2$, the last bound follows from
$(p+1)/(p-1)\le2$ and \eqref{eq:kernel-prime-power-gauss}.
For $j=1$, the bound obtained from the formula is
$(1+p^{-1/2})/(p-1)$; its ratio to $p^{-1/2}$ is
$1/(\sqrt p-1)\le2$. Thus the exponent-one case is included.

For a fixed divisor $u\mid\operatorname{rad}(q_1)$, the coprime
local mean is a product of factors from
\eqref{eq:kernel-prime-power-gauss} and
\eqref{eq:kernel-weighted-prime-power-gauss}. Its modulus is at most
$2^{\omega(u)}q_1^{-1/2}$.
The factor at primes of $D$ has modulus at most one by
\eqref{eq:kernel-root-law}. Dropping the restriction $u\le H$
after taking absolute values therefore gives
\begin{equation}
 |c(\alpha_0)|
 \le q_1^{-1/2}\sum_{u\mid\operatorname{rad}(q_1)}2^{\omega(u)}
 =3^{\omega(q_1)}q_1^{-1/2}
 \ll q_1^{-1/3}
 \le H^{-1/3}.
 \label{eq:kernel-large-coefficient}
\end{equation}
The penultimate implied constant is absolute. For example, one may use
\[
 C_*=
 \prod_{\substack{p<729\\p\text{ odd prime}}}\max(1,3p^{-1/6}),
\]
since $3\le p^{1/6}$ for every prime $p\ge729$.
Repeated prime powers only increase $q_1^{1/6}$.
This argument covers all prime powers in $q_1$, and if
$q(\alpha_D)\nmid D$, every coefficient under discussion is already
zero by \eqref{eq:kernel-root-law}.

\paragraph{Completion of the uniform comparison.}
On the minor arcs, \eqref{eq:kernel-minor-actual} and
\eqref{eq:kernel-minor-model} apply. On a major arc,
\eqref{eq:kernel-major-model}, \eqref{eq:kernel-major-actual},
\eqref{eq:kernel-small-coefficient} and
\eqref{eq:kernel-large-coefficient} apply, with $|V_N(y)|\le1$.
For clarity, every approximation error is bounded here using the fixed
value \eqref{eq:kernel-beta}:
\begin{align*}
 DHN^{-1/20}\sqrt{\log N}
   &\le N^{-0.049985}\sqrt{\log N},\\
 DH^2N^{-1/10}
   &\le N^{-0.09998},\\
 qD^2H^3/N
   &\le N^{-0.899965},\\
 \frac{qD^2H^2(1+|y|N)}{\sqrt N}
   &\le3N^{-0.29997}.
\end{align*}
In the last two lines $q\le N^{1/10}$ and
$|y|\le2N^{-9/10}$.
Each bound is $O(N^{-\beta/6})=O(H^{-1/3})$ with an absolute constant.
Together with the coefficient bound this proves
\eqref{eq:kernel-uniform-approximation} for all sufficiently large $N$.
All size and separation thresholds used above are absolute.
For the bounded remaining range, the elementary bounds
$\|S\|_\infty\ll DH$ and $\|S_0\|_\infty\le DH^2$
absorb that range by increasing the absolute constant.
Thus \eqref{eq:kernel-uniform-approximation} is uniform in all the
data of the proposition.

\paragraph{The directional bilinear identity.}
It remains to pass from the Fourier comparison to the supports in the
statement. Extend $F$ and $J$ by zero to $\mathbb Z$.
For any finitely supported kernel $h$ define
\[
 \mathcal B_h(F,J)=\frac1N\sum_{x\in\mathbb Z}
                              \sum_{n\in\mathbb Z}h(n)F(x)J(x+n).
\]
Both sums use counting measure; the displayed $1/N$ normalizes only
the base variable. Let
\[
 \widetilde F(\alpha)=\sum_{x\in\mathbb Z}F(x)e(-x\alpha),
 \qquad
 \widetilde J(\alpha)=\sum_{y\in\mathbb Z}J(y)e(-y\alpha),
 \qquad
 S_h(\alpha)=\sum_{n\in\mathbb Z}h(n)e(\alpha n).
\]
Expanding and integrating the characters imposes $y=x+n$, giving
the bilinear identity
\[
 \mathcal B_h(F,J)=\frac1N\int_0^1
 S_h(\alpha)\widetilde F(-\alpha)\widetilde J(\alpha)\,d\alpha.
\]
There are no conjugates in this identity.
Cauchy--Schwarz and counting-measure Plancherel now give
\begin{align}
 |\mathcal B_{k-k_0}(F,J)|
 &\le\frac{\|S-S_0\|_\infty}{N}
          \|F\|_{\ell^2(\mathbb Z)}\|J\|_{\ell^2(\mathbb Z)}
 \ll H^{-1/3},
 \label{eq:kernel-plancherel-error}
\end{align}
because each of the last two norms is at most $\sqrt N$.

Every term in $\mathcal B_k(F,J)$ contains $F(x)J(x+m^2)$ with
$m\ge1$, and this product is zero by hypothesis.
Thus $\mathcal B_k(F,J)=0$ term by term, regardless of the signed
weights in \eqref{eq:kernel-actual}.
For a pair $x,y$ with $F(x)J(y)\ne0$, the ordering assumption gives
$1\le y-x\le N-1$, and the residue conditions give
$y-x\equiv r\pmod D$.
Consequently both indicators in \eqref{eq:kernel-model} are automatic
for every contributing pair. Its expansion is exactly
\begin{align*}
 \mathcal B_{k_0}(F,J)
 &=\frac{D}{N^2}\sum_{x,y\in[N]}F(x)J(y)
                    \sum_{\xi\in\Xi}\mathfrak g(\xi)e(-\xi(y-x))\\
 &=D\sum_{\xi\in\Xi}\mathfrak g(\xi)
                    \widehat F(-\xi)\widehat J(\xi).
\end{align*}
Combining this with \eqref{eq:kernel-plancherel-error} proves
\eqref{eq:kernel-conclusion}.
\end{proof}

\begin{remark}
\label{rem:kernel-interval-normalization}
The normalization for later interval restrictions follows directly
from the proposition. Suppose $I_1,I_2\subseteq[N]$ are nonempty
intervals with $\max I_1<\min I_2$, and let $\Xi$ be the frequency
set in \eqref{eq:kernel-conclusion}. Suppose $F_i$ is the indicator of the restriction of a
square-difference-free set to $I_i$ and a class $a_i$ modulo $D$,
where $a_2-a_1$ is a strict square class. The coefficients
\[
 c_i(\xi)=\frac D{|I_i|}\sum_{n\in I_i}F_i(n)e(-n\xi)
          =\frac{DN}{|I_i|}\widehat F_i(\xi)
\]
satisfy
\[
 \left|\sum_{\xi\in\Xi}\mathfrak g(\xi)c_1(-\xi)c_2(\xi)\right|
 \ll H^{-1/3}\frac{DN^2}{|I_1||I_2|}.
\]
This conclusion uses the integer supports of the restrictions;
it requires no pointwise positivity of a Fourier lift.
\end{remark}
